\documentclass[10pt, compress]{beamer}

\usetheme[block=transparent, background=light]{metropolis}
\setbeamercolor{progress bar}{fg=blue3, bg=caribbeangreen}
\setbeamercolor{block title alerted}{fg=seafoam}
\setbeamercolor{frametitle}{fg=white, bg=darkgray}
\usepackage{appendixnumberbeamer}

\usepackage[export]{adjustbox}
\usepackage{booktabs}
\usepackage[scale=2]{ccicons}
\usepackage{pgfplots}
\usepgfplotslibrary{dateplot}
\usepackage{dcolumn}
\usepackage[export]{adjustbox}
\usepackage{etoolbox}

\usepackage{xspace}
\usepackage{hyperref}
\newcommand{\themename}{\textbf{\textsc{metropolis}}\xspace}

% math typesetting
\usepackage{array}
\usepackage{amsmath}
\usepackage{amssymb}
\usepackage{amsthm}
\usepackage{amsfonts}
\usepackage{relsize}
\usepackage{mathtools}
\usepackage{bm}

% tables
\usepackage{tabularx}
\usepackage{booktabs}
\usepackage{multicol}
\usepackage{multirow}

% graphics stuff
\usepackage{subfig}
\usepackage{graphicx}
\usepackage[space]{grffile} % allows us to specify directories that have spaces
\usepackage{tikz}

%\graphicspath{/Users/cassydorff/Dropbox/Research/ponyClub/JoGGSsubmission/updatedDraft/}

%CD colors
\usepackage{color}
\definecolor{lightaquablue}{rgb}{0.1,0.7,0.9}
\definecolor{seafoam}{RGB}{ 45, 200, 220}
\definecolor{deepwater}{RGB}{48, 117, 121}
\definecolor{amaranth}{rgb}{0.9, 0.17, 0.31}
\definecolor{caribbeangreen}{rgb}{0.0, 0.8, 0.6}
\definecolor{coolblack}{rgb}{0.0, 0.18, 0.39}
\definecolor{bluejazz}{rgb}{0.13, 0.67, 0.8}


% Add some colors
\definecolor{red1}{RGB}{253,219,199}
\definecolor{red2}{RGB}{244,165,130}
\definecolor{red3}{RGB}{178,24,43}

\definecolor{green1}{RGB}{229,245,224}
\definecolor{green2}{RGB}{161,217,155}
\definecolor{green3}{RGB}{49,163,84}

\definecolor{blue0}{RGB}{255,247,251}
\definecolor{blue1}{RGB}{222,235,247}
\definecolor{blue2}{RGB}{158,202,225}
\definecolor{blue3}{RGB}{49,130,189}
\definecolor{blue4}{RGB}{4,90,141}

\definecolor{purple1}{RGB}{191,211,230}
\definecolor{purple2}{RGB}{140,150,198}
\definecolor{purple3}{RGB}{140,107,177}

\definecolor{brown1}{RGB}{246,232,195}
\definecolor{brown2}{RGB}{223,194,125}
\definecolor{brown3}{RGB}{191,129,45}

% square bracket matrices
\let\bbordermatrix\bordermatrix
\patchcmd{\bbordermatrix}{8.75}{4.75}{}{}
\patchcmd{\bbordermatrix}{\left(}{\left[}{}{}
\patchcmd{\bbordermatrix}{\right)}{\right]}{}{}

% references to graphics
\makeatletter
\def\input@path{
  {/Users/janus829/Dropbox/Research/conflictEvolution/graphics/},   {/Users/s7m/Dropbox/Research/conflictEvolution/graphics/},  {/Users/cassydorff/Dropbox/Research/nothingbutnet/conflictEvolution/graphics/}
     % {/Users/maxgallop/Dropbox (Netgroup)/conflictEvolutionMex/graphics/} % paren in path screws stuff up
}

\graphicspath{
  {/Users/janus829/Dropbox/Research/conflictEvolution/graphics/},   {/Users/s7m/Dropbox/Research/conflictEvolution/graphics/},  {/Users/cassydorff/Dropbox/Research/nothingbutnet/conflictEvolution/graphics/}
      % {/Users/maxgallop/Dropbox (Netgroup)/conflictEvolutionMex/graphics/}  % paren in path screws stuff up
}

\title[Nigeria Conflict]{\textsc{Predicting Violence: Network Dynamics in Nigeria }}
\author[Dorff, Gallop, \& Minhas]{Cassy Dorff (University of New Mexico),\\ Max Gallop (University of Strathclyde),\\ and Shahryar Minhas (Michigan State University)\\} 

\date{\textsc{\today}}

\begin{document}
%\begin{frame}[plain]
%	
%	Me
%	
%\end{frame}

\frame{\titlepage}

\section{Motivation}
%%%%%%%%%%%%%%%%%%%%%%%%%%%%%%%%%%%%%%%%%%%%%%%%%%%%%%%%%%%%

%%%%%%%%%%%%%%%%%%%%%%%%%%%%%%%%%%%%%%%%%%%%%%%%%%%%%%%%%%%%

\frame
{
  \frametitle{Intrastate War}
  Extensive literature on the causes and prediction of intrastate conflict
  \begin{table}
  \begin{tabular}{lc}
    Hegre et al. (2001) & \multirow{5}{*}{\includegraphics[width=.6\textwidth]{confGraphs.png}} \\
    Fearon \& Laitin (2003) & \\
    Collier et al. (2004) & \\
    Salehyan (2013)  & \\
    K.G. Cunningham (2013) & \\
    Sambanis \& Shayo (2013) & \\
    Lacina (2014) & \\
    Prorok (2016) & \\
  \end{tabular}
  \end{table}

} 
%%%%%%%%%%%%%%%%%%%%%%%%%%%%%%%%%%%%%%%%%%%%%%%%%%%%%%%%%%%%


%%%%%%%%%%%%%%%%%%%%%%%%%%%%%%%%%%%%%%%%%%%%%%%%%%%%%%%%%%%%

\frame
{
  \frametitle{Intrastate War}
  Extensive literature on the causes and prediction of intrastate conflict
  \begin{table}
  \begin{tabular}{lc}
    Hegre et al. (2001) & \multirow{5}{*}{\includegraphics[width=.6\textwidth]{confGraphs.png}} \\
    Fearon \& Laitin (2003) & \\
    Collier et al. (2004) & \\
    Salehyan (2013)  & \\
    K.G. Cunningham (2013) & \\
    Sambanis \& Shayo (2013) & \\
    Lacina (2014) & \\
    Prorok (2016) & \\
  \end{tabular}
  \end{table}

\vspace{.2in}
 \textbf{\textcolor{purple3}{Fearon \& Laitin (2003) has been cited over 6,000 times!}}

} 

%%%%%%%%%%%%%%%%%%%%%%%%%%%%%%%%%%%%%%%%%%%%%%%%%%%%%%%%%%%%

\frame{
  \frametitle{Conflicts are complex: unpacking social structure}

  Roughly a \textbf{third} of all intrastate conflict between 1989 and 2003 have been fought with multiple warring parties (UCDP/PRIO 2007).
  \vspace{.1in}
  
 Conflicts involve multiple actors with changing relationships overtime
 
    \vspace{.3in}

}


%%%%%%%%%%%%%%%%%%%%%%%%%%%%%%%%%%%%%%%%%%%%%%%%%%%%%%%%%%%%
\frame{
  \frametitle{Conflicts are complex: unpacking social structure}

  Roughly a \textbf{third} of all intrastate conflict between 1989 and 2003 have been fought with multiple warring parties (UCDP/PRIO 2007).
  \vspace{.1in}
  
 Conflicts involve multiple actors with changing relationships overtime
 
 \begin{itemize}
 \item Coordination (Bakke et al 2012; Findley \& Rudloff, 2012)
 \item Spoiler groups and veto-players (Cunningham, 2006)
 \item Disaggregating actors (Shellman et al, 2010)
\end{itemize}
}

%%%%%%%%%%%%%%%%%%%%%%%%%%%%%%%%%%%%%%%%%%%%%%%%%%%%%%%%%%%%

\frame{
  \frametitle{Pairing Empirical Analysis to Theory}

  \begin{quote}
 ``Existence of \textcolor{amaranth}{multiple rebel groups} means we can no longer understand civil wars with a sole focus on state attributes. In fact, the government's strategies leading to victory, defeat, or continuation of war can only be understood \textcolor{amaranth}{in relation to} the rebel group/groups it is fighting." 
  \end{quote}
  \vspace{-.1in}
 \center  Akcinaroglu (2012)
  
}
%%%%%%%%%%%%%%%%%%%%%%%%%%%%%%%%%%%%%%%%%%%%%%%%%%%%%%%%%%%%
%%%%%%%%%%%%%%%%%%%%%%%%%%%%%%%%%%%%%%%%%%%%%%%%%%%%%%%%%%%%
\frame{
	\frametitle{Our Project}
	\textbf{Conflict processes are driven by the evolution of relationships overtime. }
		\begin{enumerate}
			\item Intrastate conflicts $\rightarrow$ single complex system composed of multiple actors in conflict 
			\end{enumerate}  
\vspace{1in}
}

\frame{
	\frametitle{Our Project}
	\textbf{Conflict processes are driven by the evolution of relationships overtime. }
		\begin{enumerate}
			\item Intrastate conflicts $\rightarrow$ single complex system composed of multiple actors in conflict 
			\item Armed actors \& battles $=$ nodes and ties in a network
		\end{enumerate}  
\vspace{.7in}
}

\frame{
	\frametitle{Our Project}
	\textbf{Conflict processes are driven by the evolution of relationships overtime. }
		\begin{enumerate}
			\item Intrastate conflicts $\rightarrow$ single complex system composed of multiple actors in conflict 
			\item Armed actors \& battles $=$ nodes and ties in a network
			\item Novel model captures interdependencies across actors within the conflict system
		\end{enumerate}  

}

\frame{
	\frametitle{Our Project}
	\textbf{Conflict processes are driven by the evolution of relationships overtime. }
		\begin{enumerate}
			\item Intrastate conflicts $\rightarrow$ single complex system composed of multiple actors in conflict 
			\item Armed actors \& battles $=$ nodes and ties in a network
			\item Novel model captures relationships endogenous to the conflict system 
			\item Our approach provides unbiased parameter estimates \& out performs standard approaches
		\end{enumerate}  

}

\frame{
	\frametitle{Our Project}
	\textbf{Conflict processes are driven by the evolution of relationships overtime. }
		\begin{enumerate}
			\item Intrastate conflicts $\rightarrow$ single complex system composed of multiple actors in conflict 
			\item Armed actors \& battles $=$ nodes and ties in a network
			\item Novel model captures relationships endogenous to the conflict system 
			\item Our approach provides precise estimates, \& out performs standard approaches
			\item Uncovers important relational patterns of conflict with substantive implications for the study of conflict processes
		\end{enumerate}  

}


%%%%%%%%%%%%%%%%%%%%%%%%%%%%%%%%%%%%%%%%%%%%%%%%%%%%%%%%%%%%
%%%%%%%%%%%%%%%%%%%%%%%%%%%%%%%%%%%%%%%%%%%%%%%%%%%%%%%%%%%%

\section{Networks \& Conflict Processes}

\frame{
  \frametitle{From dyads to networks}

  Dyadic data consists of a set of:
  \begin{itemize}
    \item nodes (e.g., rebel group actors)
    \item measurements specific to a pair of actors (e.g., the occurrence of a battle)
  \end{itemize}

  \centering
  \includegraphics[width=.9\textwidth]{df_adj_net3.png}

}
%%%%%%%%%%%%%%%%%%%%%%%%%%%%%%%%%%%%%%%%%%%%%%%%%%%%%%%%%%%%
%%%%%%%%%%%%%%%%%%%%%%%%%%%%%%%%%%%%%%%%%%%%%%%%%%%%%%%%%%%%

%\frame{
%  \frametitle{Dyadic data assumptions}
%
%GLM: $y_{ij} \sim \beta^{T} X_{ij} + e_{ij}$
%
%Networks typically show evidence against independence of dyadic interactions
%
%Not accounting for dependence can lead to:
%
%\begin{itemize}
%\item biased effects estimation
%\item uncalibrated confidence intervals
%\item poor predictive performance
%\item inaccurate description of network phenomena
%\end{itemize}
%
%}

%%%%%%%%%%%%%%%%%%%%%%%%%%%%%%%%%%%%%%%%%%%%%%%%%%%%%%%%%%%%
%%%%%%%%%%%%%%%%%%%%%%%%%%%%%%%%%%%%%%%%%%%%%%%%%%%%%%%%%%%%

\frame{
  \frametitle{Network Effects}
 How does evolution in the structure of relationships influence conflict over time?
\begin{itemize}
\item 1st-order: Sender effects\\
\item 2nd-order: Reciprocity\\
\item 3rd-order: Homophily \& Stochastic equivalence\\
\item System level: Changing actor composition\\
\end{itemize}
}

%%%%%%%%%%%%%%%%%%%%%%%%%%%%%%%%%%%%%%%%%%%%%%%%%%%%%%%%%%%%
%%%%%%%%%%%%%%%%%%%%%%%%%%%%%%%%%%%%%%%%%%%%%%%%%%%%%%%%%%%%
\frame{
  \frametitle{Network phenomena: sender heterogeneity}

  Values across a row, say $\{y_{ij},y_{ik},y_{il}\}$, may be more similar to each other than other values in the adjacency matrix because each of these values has a common sender $i$

  \centering
  \includegraphics[width=.7\textwidth]{adjRowDep}

}

%%%%%%%%%%%%%%%%%%%%%%%%%%%%%%%%%%%%%%%%%%%%%%%%%%%%%%%%%%%%
%%%%%%%%%%%%%%%%%%%%%%%%%%%%%%%%%%%%%%%%%%%%%%%%%%%%%%%%%%%%

\frame{
  \frametitle{Network phenomena: receiver heterogeneity}

  Values across a column, say $\{y_{ji},y_{ki},y_{li}\}$, may be more similar to each other than other values in the adjacency matrix because each of these values has a common receiver $i$

  \centering
  \includegraphics[width=.7\textwidth]{adjColDep}

}
%%%%%%%%%%%%%%%%%%%%%%%%%%%%%%%%%%%%%%%%%%%%%%%%%%%%%%%%%%%%

%%%%%%%%%%%%%%%%%%%%%%%%%%%%%%%%%%%%%%%%%%%%%%%%%%%%%%%%%%%%
\frame{
  \frametitle{Network phenomena: sender-receiver covariance}

  Actors who are more likely to send ties in a network may also be more likely to receive them

  \centering
  \includegraphics[width=.7\textwidth]{adjRowColCovar}

}
%%%%%%%%%%%%%%%%%%%%%%%%%%%%%%%%%%%%%%%%%%%%%%%%%%%%%%%%%%%%

%%%%%%%%%%%%%%%%%%%%%%%%%%%%%%%%%%%%%%%%%%%%%%%%%%%%%%%%%%%%
\frame{
  \frametitle{Network phenomena: reciprocity}

  Values of $y_{ij}$ and $y_{ji}$ may be statistically dependent

  \centering
  \includegraphics[width=.7\textwidth]{adjRecip}

}
%%%%%%%%%%%%%%%%%%%%%%%%%%%%%%%%%%%%%%%%%%%%%%%%%%%%%%%%%%%%

%%%%%%%%%%%%%%%%%%%%%%%%%%%%%%%%%%%%%%%%%%%%%%%%%%%%%%%%%%%%
\frame{
  \frametitle{Network phenomena: third order dependencies}

  \vspace{-10mm}

  \begin{table}[ht]
  \begin{tabular}{lcr}
  \qquad\scshape{Homophily} & & \scshape{Stochastic Equivalence} \\
  \includegraphics[width=.33\textwidth]{homophNet.pdf} & \hspace{2cm} &
  \includegraphics[width=.33\textwidth]{stochEquivNet.pdf}  
  \end{tabular}
  \end{table}
  }
%%%%%%%%%%%%%%%%%%%%%%%%%%%%%%%%%%%%%%%%%%%%%%%%%%%%%%%%%%%%


%%%%%%%%%%%%%%%%%%%%%%%%%%%%%%%%%%%%%%%%%%%%%%%%%%%%%%%%%%%%

\frame{
  \frametitle{Network phenomena: changing actor composition}
\begin{center}
  \includegraphics[width=1\textwidth]{netsOverTmebw}
\end{center}

}
%%%%%%%%%%%%%%%%%%%%%%%%%%%%%%%%%%%%%%%%%%%%%%%%%%%%%%%%%%%%

%  To account for these patterns we can build on the SRM framework and find an expression for $\gamma$:
%  
%  \vspace{-5mm}
%  \begin{align*}
%  \centering
%  y_{ij} &\approx \beta^{T} X_{ij} + a_{i} + b_{j} + \gamma(u_{i},v_{j})
%  \end{align*}
%}

%%%%%%%%%%%%%%%%%%%%%%%%%%%%%%%%%%%%%%%%%%%%%%%%%%%%%%%%%%%%
\section{Nigeria}

\frame{
	\frametitle{Intrastate Conflict: Nigeria's intrastate conflict system}
Complex, multi-actor conflict
\begin{itemize}
\item  \textcolor{bluejazz}{Numerous} violent political groups including ethnic militias, militant regional groups and Islamist insurgents
\item Political violence of \textcolor{bluejazz}{all types} has risen substantially since 2011 with violence against civilians seeing the most dramatic increase.
\item \textcolor{bluejazz}{Civilians} engage in both violent and nonviolent resistance efforts.
\end{itemize}


}

\frame{
	\frametitle{Intrastate Conflict: Nigeria's intrastate conflict system}
Complex, multi-actor conflict
\begin{itemize}
\item  \textcolor{bluejazz}{numerous} violent political groups including ethnic militias, militant regional groups and Islamist insurgents
\item political violence of \textcolor{bluejazz}{all types} has risen substantially since 2011 with violence against civilians seeing the most dramatic increase.
\item \textcolor{bluejazz}{civilians} engage in both violent and nonviolent resistance efforts.
\end{itemize}
  \includegraphics[width=1\textwidth]{actorOvertime}

}

\frame{
  \frametitle{Intrastate Conflict: Nigeria's intrastate conflict system}
  \begin{center}
  \includegraphics[width=1\textwidth]{nigeria_2000_2016.pdf}
  \end{center}

}
%%%%%%%%%%%%%%%%%%%%%%%%%%%%%%%%%%%%%%%%%%%%%%%%%%%%%%%%%%%%

%%%%%%%%%%%%%%%%%%%%%%%%%%%%%%%%%%%%%%%%%%%%%%%%%%%%%%%%%%%%
\frame{
  \frametitle{Spatial Distribution of Conflict Pre Boko Haram}

%  \vspace{-5mm}
  \centering
  \includegraphics[width=1\textwidth]{nigConfMap_preBH.pdf}

}
%%%%%%%%%%%%%%%%%%%%%%%%%%%%%%%%%%%%%%%%%%%%%%%%%%%%%%%%%%%%

%%%%%%%%%%%%%%%%%%%%%%%%%%%%%%%%%%%%%%%%%%%%%%%%%%%%%%%%%%%%
\frame{
  \frametitle{Spatial Distribution of Conflict Post Boko Haram}

  \centering
  \includegraphics[width=1\textwidth]{nigConfMap_postBH.pdf}

}
%%%%%%%%%%%%%%%%%%%%%%%%%%%%%%%%%%%%%%%%%%%%%%%%%%%%%%%%%%%%

%%%%%%%%%%%%%%%%%%%%%%%%%%%%%%%%%%%%%%%%%%%%%%%%%%%%%%%%%%%%
%\frame{
%  \frametitle{Boko Haram's Entrance in Network}
%
%  \centering
%  \includegraphics[width=1\textwidth]{nigeria_postBK}
%
%}
%%%%%%%%%%%%%%%%%%%%%%%%%%%%%%%%%%%%%%%%%%%%%%%%%%%%%%%%%%%%


%%%%%%%%%%%%%%%%%%%%%%%%%%%%%%%%%%%%%%%%%%%%%%%%%%%%%%%%%%%%
\section{Modeling Approach \& Results}
\frame{
  \frametitle{Social Relations Model (The ``A'' in AME)}

Additive effects portion of AME (Warner et al. 1979; Li \& Loken 2002):

\begin{align*}
\begin{aligned}
      y_{ij} &= \color{amaranth}{\mu} + \color{amaranth}{e_{ij}} \\
      e_{ij} &= a_{i} + b_{j} + \epsilon_{ij} \\
      \{ (a_{1}, b_{1}), \ldots, (a_{n}, b_{n}) \} &\sim N(0,\Sigma_{ab}) \\ 
      \{ (\epsilon_{ij}, \epsilon_{ji}) : \; i \neq j\} &\sim N(0,\Sigma_{\epsilon}), \text{ where } \\
      \Sigma_{ab} = \begin{pmatrix} \sigma_{a}^{2} & \sigma_{ab} \\ \sigma_{ab} & \sigma_{b}^2   \end{pmatrix} \;\;\;\;\; &\Sigma_{\epsilon} = \sigma_{\epsilon}^{2} \begin{pmatrix} 1 & \rho \\ \rho & 1  \end{pmatrix}
\end{aligned}
\end{align*}

\begin{itemize}
\item $\mu$ baseline measure of network activity (for the purpose of regression we turn this into $\bm\beta^{T} \mathbf{X}_{ij,t}$)
\item $e_{ij}$ residual variation that we will use the SRM to decompose
\end{itemize}

}
%%%%%%%%%%%%%%%%%%%%%%%%%%%%%%%%%%%%%%%%%%%%%%%%%%%%%%%%%%%%

%%%%%%%%%%%%%%%%%%%%%%%%%%%%%%%%%%%%%%%%%%%%%%%%%%%%%%%%%%%%


\frame{
  \frametitle{Social Relations Model (The ``A'' in AME)}

\begin{align*}
\begin{aligned}
      y_{ij} &= \mu + e_{ij} \\
      e_{ij} &= \color{amaranth}{a_{i} + b_{j}} + \epsilon_{ij} \\
      \color{amaranth}{\{ (a_{1}, b_{1}), \ldots, (a_{n}, b_{n}) \}} &\sim N(0,\Sigma_{ab}) \\ 
      \{ (\epsilon_{ij}, \epsilon_{ji}) : \; i \neq j\} &\sim N(0,\Sigma_{\epsilon}), \text{ where } \\
      \Sigma_{ab} = \begin{pmatrix} \sigma_{a}^{2} & \sigma_{ab} \\ \sigma_{ab} & \sigma_{b}^2   \end{pmatrix} \;\;\;\;\; &\Sigma_{\epsilon} = \sigma_{\epsilon}^{2} \begin{pmatrix} 1 & \rho \\ \rho & 1  \end{pmatrix}
\end{aligned}
\end{align*}

\begin{itemize}
\item row/sender effect ($a_{i}$) \& column/receiver effect ($b_{j}$)
\item Modeled jointly to account for correlation in how active an actor is in sending and receiving ties
\end{itemize}

}
%%%%%%%%%%%%%%%%%%%%%%%%%%%%%%%%%%%%%%%%%%%%%%%%%%%%%%%%%%%%

%%%%%%%%%%%%%%%%%%%%%%%%%%%%%%%%%%%%%%%%%%%%%%%%%%%%%%%%%%%%
\frame{
  \frametitle{Social Relations Model (The ``A'' in AME)}

\begin{align*}
\begin{aligned}
      y_{ij} &= \mu + e_{ij} \\
      e_{ij} &= a_{i} + b_{j} + \epsilon_{ij} \\
      \{ (a_{1}, b_{1}), \ldots, (a_{n}, b_{n}) \} &\sim N(0,\color{amaranth}{\Sigma_{ab}}) \\ 
      \{ (\epsilon_{ij}, \epsilon_{ji}) : \; i \neq j\} &\sim N(0,\Sigma_{\epsilon}), \text{ where } \\
      \color{amaranth}{\Sigma_{ab}} = \begin{pmatrix} \sigma_{a}^{2} & \sigma_{ab} \\ \sigma_{ab} & \sigma_{b}^2   \end{pmatrix} \;\;\;\;\; &\Sigma_{\epsilon} = \sigma_{\epsilon}^{2} \begin{pmatrix} 1 & \rho \\ \rho & 1  \end{pmatrix}
\end{aligned}
\end{align*}

\begin{itemize}
\item $\sigma_{a}^{2}$ and $\sigma_{b}^{2}$ capture heterogeneity in the row and column means
\item $\sigma_{ab}$ describes the linear relationship between these two effects (i.e., whether actors who send [receive] a lot of ties also receive [send] a lot of ties)
\end{itemize}

}
%%%%%%%%%%%%%%%%%%%%%%%%%%%%%%%%%%%%%%%%%%%%%%%%%%%%%%%%%%%%

%%%%%%%%%%%%%%%%%%%%%%%%%%%%%%%%%%%%%%%%%%%%%%%%%%%%%%%%%%%%
\frame{
  \frametitle{Social Relations Model (The ``A'' in AME)}

\begin{align*}
\begin{aligned}
      y_{ij} &= \mu + e_{ij} \\
      e_{ij} &= a_{i} + b_{j} + \color{amaranth}{\epsilon_{ij}} \\
      \{ (a_{1}, b_{1}), \ldots, (a_{n}, b_{n}) \} &\sim N(0,\Sigma_{ab}) \\ 
      \color{amaranth}{\{ (\epsilon_{ij}, \epsilon_{ji}) : \; i \neq j\}} &\sim N(0,\color{amaranth}{\Sigma_{\epsilon}}), \text{ where } \\
      \Sigma_{ab} = \begin{pmatrix} \sigma_{a}^{2} & \sigma_{ab} \\ \sigma_{ab} & \sigma_{b}^2   \end{pmatrix} \;\;\;\;\; & \color{red}{\Sigma_{\epsilon}} = \sigma_{\epsilon}^{2} \begin{pmatrix} 1 & \rho \\ \rho & 1  \end{pmatrix}
\end{aligned}
\end{align*}

\begin{itemize}
\item $\epsilon_{ij}$ captures the within dyad effect
\item Second-order dependencies are described by $\sigma_{\epsilon}^{2}$
\item Reciprocity, aka within dyad correlation, represented by $\rho$
\end{itemize}
}
%%%%%%%%%%%%%%%%%%%%%%%%%%%%%%%%%%%%%%%%%%%%%%%%%%%%%%%%%%%%

%%%%%%%%%%%%%%%%%%%%%%%%%%%%%%%%%%%%%%%%%%%%%%%%%%%%%%%%%%%%
\frame{
  \frametitle{Latent Factor Model: The ``M'' in AME}

Each node $i$ has an unknown latent factor

\vspace{-5mm}
\begin{align*}
{\textbf{u}_{i},\textbf{v}_{j}} \in \mathbb{R}^{k} \;\; i,j \in \{1, \ldots, n \} \\
\end{align*}

\vspace{-5mm}
The probability of a tie from $i$ to $j$ depends on their latent factors

\vspace{-5mm}
\begin{align*}
\begin{aligned}
  \gamma(\textbf{u}_{i}, \textbf{v}_{j}) &= \textbf{u}_{i}^{T} D \textbf{v}_{j} \\
  &= \sum_{k \in K} d_{k} u_{ik} v_{jk} \\
  &D \text{ is a  } K \times K \text{ diagonal matrix}
\end{aligned}
\end{align*}

Accounts for both stochastic equivalence and homophily (Hoff 2008)

}
%%%%%%%%%%%%%%%%%%%%%%%%%%%%%%%%%%%%%%%%%%%%%%%%%%%%%%%%%%%%

%%%%%%%%%%%%%%%%%%%%%%%%%%%%%%%%%%%%%%%%%%%%%%%%%%%%%%%%%%%%
\frame{
  \frametitle{Additive and Multiplicative Effects (AME) Model}

  \begin{align*}
    \begin{aligned}
      y_{ij,t} &= g(\theta_{ij,t}) \\ 
      &\theta_{ij,t} = \bm\beta^{T} \mathbf{X}_{ij,t} + e_{ij,t} \\
      % &\theta_{ij,t} = \bm\beta_{d}^{T} \mathbf{X}_{ij,t} + \bm\beta_{s}^{T} \mathbf{X}_{i,t} + \bm\beta_{r}^{T} \mathbf{X}_{j,t} + e_{ij,t} \\      
      &e_{ij,t} = a_{i} + b_{j}  + \epsilon_{ij} + \alpha(\textbf{u}_{i}, \textbf{v}_{j}) \text{  , where } \\
      &\qquad \alpha(\textbf{u}_{i}, \textbf{v}_{j}) = \textbf{u}_{i}^{T} \textbf{D} \textbf{v}_{j} = \sum_{k \in K} d_{k} u_{ik} v_{jk} \\ 
    \end{aligned}
  \end{align*}

  (Hoff 2005; Hoff 2008; Hoff et al. 2013; Minhas et al. 2016)

}
%%%%%%%%%%%%%%%%%%%%%%%%%%%%%%%%%%%%%%%%%%%%%%%%%%%%%%%%%%%%

%%%%%%%%%%%%%%%%%%%%%%%%%%%%%%%%%%%%%%%%%%%%%%%%%%%%%%%%%%%%

%%%%%%%%%%%%%%%%%%%%%%%%%%%%%%%%%%%%%%%%%%%%%%%%%%%%%%%%%%%%
\frame{
  \frametitle{Data}

  Armed Conflict Location and Event Data Project (ACLED) developed by Raleigh et al. (2010)

  \begin{itemize}
    \item ACLED records armed conflict and protest events in over 60 developing countries
    \item We use ACLED \textit{battles} data for Nigeria to generate a measure of conflict where:
    \begin{itemize}
      \item $y_{ij,t} = 1$ indicates that a conflict occurred when actor $i$ attacked actor $j$ at time $t$
      \item $y_{ij,t} = 0$ if no conflict occurred
    \end{itemize}
    \item We focus only on modeling the interactions between armed groups that are engaged in battles for at least 5 years during the 2000-2016 period, which results in a total of 37 armed groups
  \end{itemize}

}
%%%%%%%%%%%%%%%%%%%%%%%%%%%%%%%%%%%%%%%%%%%%%%%%%%%%%%%%%%%%

%%%%%%%%%%%%%%%%%%%%%%%%%%%%%%%%%%%%%%%%%%%%%%%%%%%%%%%%%%%%
\frame{
  \frametitle{Covariates}

  \begin{itemize}
    \item Country-Level covariates:
    \begin{itemize}
      \item Post Boko-Haram
      \item Neighborhood conflict
      \item Election year      
    \end{itemize}
    \item Sender and Receiver-Level Covariates:
    \begin{itemize}
      \item Violence against civilians
      \item Riots/Protests directed against actor
      \item Geographic spread
    \end{itemize}
  \end{itemize}

}
%%%%%%%%%%%%%%%%%%%%%%%%%%%%%%%%%%%%%%%%%%%%%%%%%%%%%%%%%%%%

%%%%%%%%%%%%%%%%%%%%%%%%%%%%%%%%%%%%%%%%%%%%%%%%%%%%%%%%%%%%
\frame{
  \frametitle{Model Results}

  \centering
  \begin{tabular}{c}
  \includegraphics[height=.7\textheight]{betaEst} \\
  \includegraphics[height=.2\textheight]{vcEst} \\  
  \end{tabular}

}
%%%%%%%%%%%%%%%%%%%%%%%%%%%%%%%%%%%%%%%%%%%%%%%%%%%%%%%%%%%%

%%%%%%%%%%%%%%%%%%%%%%%%%%%%%%%%%%%%%%%%%%%%%%%%%%%%%%%%%%%%
\frame{
  \frametitle{Multiplicative Effects}

  \centering
  \includegraphics[width=1\textwidth]{circPlotSimplev2}  

}
%%%%%%%%%%%%%%%%%%%%%%%%%%%%%%%%%%%%%%%%%%%%%%%%%%%%%%%%%%%%

%%%%%%%%%%%%%%%%%%%%%%%%%%%%%%%%%%%%%%%%%%%%%%%%%%%%%%%%%%%%
\frame{
  \frametitle{Out of Sample Cross-Validation}

  \centering
  \begin{tabular}{ll}
  \includegraphics[width=.5\textwidth]{roc_outSample} & 
  \includegraphics[width=.5\textwidth]{rocPr_outSample}
  \end{tabular}  

}
%%%%%%%%%%%%%%%%%%%%%%%%%%%%%%%%%%%%%%%%%%%%%%%%%%%%%%%%%%%%

%%%%%%%%%%%%%%%%%%%%%%%%%%%%%%%%%%%%%%%%%%%%%%%%%%%%%%%%%%%%
\frame{
  \frametitle{Out of Sample Forecast}

  \centering
  \begin{tabular}{ll}
  \includegraphics[width=.5\textwidth]{roc_outSampleForecast_1} & 
  \includegraphics[width=.5\textwidth]{rocPr_outSampleForecast_1}
  \end{tabular}  

}
%%%%%%%%%%%%%%%%%%%%%%%%%%%%%%%%%%%%%%%%%%%%%%%%%%%%%%%%%%%%

% %%%%%%%%%%%%%%%%%%%%%%%%%%%%%%%%%%%%%%%%%%%%%%%%%%%%%%%%%%%%
 \frame{
   \frametitle{Future work}
   Are \textcolor{blue3}{``people-power"} movements less effective in multi-actor civil conflicts?
\vspace{.08in}

Why \textcolor{blue3}{does violence against civilians} increase an actor's conflictual behavior towards armed groups?
\vspace{.08in}

Does our \textcolor{blue3}{``key player"} effect matter in other conflict settings?
}

\frame{
\frametitle{Key Take-aways}
  \textbf{\textcolor{seafoam}{\textsc{confirmed}}}: Intrastate conflict is a network process! Structure of relationships influences violence between actors (reciprocity and warring communities characterize social patterns in the data) 
  \vspace{.08in}
    
   \textbf{\textcolor{seafoam}{\textsc{confirmed}}}: Key players  alter violence in the conflict system, even in warring dyads the key player is not directly involved in.
  \vspace{.08in}
  
   \textbf{\textcolor{seafoam}{\textsc{confirmed}}}: Network model of conflict out performs standard approaches
  

}
% %%%%%%%%%%%%%%%%%%%%%%%%%%%%%%%%%%%%%%%%%%%%%%%%%%%%%%%%%%%%

%%%%%%%%%%%%%%%%%%%%%%%%%%%%%%%%%%%%%%%%%%%%%%%%%%%%%%%%%%%%
\plain{Thanks!\\
\small
\textsc{cassydorff.com}}
%%%%%%%%%%%%%%%%%%%%%%%%%%%%%%%%%%%%%%%%%%%%%%%%%%%%%%%%%%%%

%%%%%%%%%%%%%%%%%%%%%%%%%%%%%%%%%%%%%%%%%%%%%%%%%%%%%%%%%%%%
\frame{
  \frametitle{Network GOF}

  \centering
  \includegraphics[width=.8\textwidth]{netGOF}  

}
%%%%%%%%%%%%%%%%%%%%%%%%%%%%%%%%%%%%%%%%%%%%%%%%%%%%%%%%%%%%
%%%%%%%%%%%%%%%%%%%%%%%%%%%%%%%%%%%%%%%%%%%%%%%%%%%%%%%%%%%%
\frame{
  \frametitle{Additive Sender/Receiver Random Effects}
  \centering
  \includegraphics[width=.9\textwidth]{abEst2}

}
%%%%%%%%%%%%%%%%%%%%%%%%%%%%%%%%%%%%%%%%%%%%%%%%%%%%%%%%%%%%

%%%%%%%%%%%%%%%%%%%%%%%%%%%%%%%%%%%%%%%%%%%%%%%%%%%%%%%%%%%%

\frame{
  \frametitle{Dyadic data assumptions}

GLM: $y_{ij} \sim \beta^{T} X_{ij} + e_{ij}$

Networks typically show evidence against independence of dyadic interactions

Not accounting for dependence can lead to:

\begin{itemize}
\item biased effects estimation
\item uncalibrated confidence intervals
\item poor predictive performance
\item inaccurate description of network phenomena
\end{itemize}

We've been hearing this concern for decades now:

\begin{tabular}{lll}
Thompson \& Walker (1982) & Beck et al. (1998) & Snijders (2011) \\
Frank \& Strauss (1986) & Signorino (1999) & Erikson et al. (2014) \\
Kenny (1996) & Li \& Loken (2002) & Aronow et al. (2015) \\
Krackhardt (1998) & Hoa \& Ward (2004) & Athey et al. (2016) \\
\end{tabular}

}

\frame{
	\frametitle{ACLED Data - Nigeria}
	Data collection
	\begin{itemize}
	\item Battles are violent clashes between at least two armed groups. 
	\item Battles make up approximately one third of the dataset.
	\item Data types: civic society (reports, NGOs), media (newspapers), Analysts (specialists' reports), governing bodies (UN reports), ``Local source project" (ACLED is connected with local sources)
	\item Analysis of data does not reveal urban bias
	\end{itemize} 
	
	
}

%%%%%%%%%%%%%%%%%%%%%%%%%%%%%%%%%%%%%%%%%%%%%%%%%%%%%%%%%%%%
\frame{
  \frametitle{Boko Haram's Entrance in Network}
  \centering
  \includegraphics[width=1\textwidth]{nigeria_postBK}

}

\frame{
  \frametitle{ERGMs}
ERGMs are useful when researchers are interested in the role that a specific list of network statistics have in giving rise to a certian network. (Such as: number of transitive triads in a network, balanced triads, reciprocal pairs, etc.)
\begin{itemize}
	\item ERGMs provide a way to find the probability of a network given the patterns it exhibits
	\item the researcher must specify which network statistics should give rise to a particular network of interest
\end{itemize}

}
%%%%%%%%%%%%%%%%%%%%%%%%%%%%%%%%%%%%%%%%%%%%%%%%%%%%%%%%%%%%



\end{document}
